\section{Data, Probability \& Descriptive Statistics}

\begin{frame}{Learning Goals}
\begin{itemize}
  \item Distinguish data types (categorical vs. numerical, discrete vs. continuous) and read frequency distributions, histograms, and box plots.
  \pause
  \item Compute and interpret measures of center (mean, median, mode) and spread (variance, standard deviation, IQR).
  \pause
  \item Define sample spaces and events, and apply the complement, addition, and multiplication rules of probability.
  \pause
  \item Distinguish mutual exclusivity from independence, and empirical from classical probability.
\end{itemize}
\end{frame}

\begin{frame}{Reading Resources}
\begin{itemize}
  \item Primary reference for Module 0, Lectures 1--5: Illowsky \& Dean, \textit{Introductory Statistics 2e} (OpenStax), free at \url{https://openstax.org/details/books/introductory-statistics-2e}.
  \pause
  \item We omit topics not critical for AI: chi-square tests, two-sample $t$-tests, ANOVA, and study-design methodology.
  \pause
  \item Hypothesis testing is deferred to Lecture 10 (Statistical Model Evaluation), its most relevant context for this course.
\end{itemize}
\end{frame}

\begin{frame}{Why Statistics Before Machine Learning}
\begin{itemize}
  \item Every ML model is trained on a \textbf{sample} and evaluated by \textbf{estimating} how well it will generalize.
  \pause
  \item Loss functions (cross-entropy, MSE) are maximum-likelihood objectives in disguise --- Lecture 5.
  \pause
  \item Bayes' theorem underlies methods from Naive Bayes to diffusion models --- Lecture 2.
  \pause
  \item This module builds the probability and estimation vocabulary that Modules 1--3 assume.
\end{itemize}
\end{frame}

\begin{frame}{Data Types}
\begin{columns}[T]
\column{0.48\textwidth}
\textbf{Categorical (qualitative)}
\begin{itemize}
  \item nominal: no natural order (e.g. color, model architecture name)
  \pause
  \item ordinal: ordered, unequal spacing (e.g. rating: low/medium/high)
\end{itemize}

\pause
\column{0.48\textwidth}
\textbf{Numerical (quantitative)}
\begin{itemize}
  \item discrete: countable (e.g. number of tokens)
  \pause
  \item continuous: measured on a scale (e.g. inference latency)
\end{itemize}
\end{columns}

\pause
\begin{rlintuition}
The data type determines which summary statistics and plots are valid: means and standard deviations only make sense for numerical data; categorical data is summarized with counts and proportions.
\end{rlintuition}
\end{frame}

\begin{frame}{Measures of Center}
\begin{rldefinition}
For a sample $x_1, \dots, x_n$:
\[
\bar{x} = \frac{1}{n}\sum_{i=1}^{n} x_i
\qquad
\text{median} = \text{middle value of sorted } x_i
\]
\end{rldefinition}
\pause
\begin{itemize}
  \item The mean is sensitive to outliers; the median is robust to them.
  \pause
  \item The mode is the most frequent value --- the only measure of center defined for categorical data.
  \pause
  \item Skewed distributions (e.g. income, latency) are usually better summarized by the median.
\end{itemize}
\end{frame}

\begin{frame}{Measures of Spread}
\begin{rldefinition}
\[
s^2 = \frac{1}{n-1}\sum_{i=1}^{n}(x_i - \bar{x})^2
\qquad
s = \sqrt{s^2}
\]
\end{rldefinition}
\pause
\begin{itemize}
  \item $s^2$ is the sample variance; dividing by $n-1$ (Bessel's correction) gives an unbiased estimator of the population variance.
  \pause
  \item The interquartile range $\text{IQR} = Q_3 - Q_1$ is a robust spread measure, used to flag outliers ($< Q_1 - 1.5\,\text{IQR}$ or $> Q_3 + 1.5\,\text{IQR}$).
  \pause
  \item Box plots visualize $Q_1$, median, $Q_3$, and outliers in one figure.
\end{itemize}
\end{frame}

\begin{frame}{Sample Spaces, Events \& Probability Rules}
\begin{rldefinition}
The \textbf{sample space} $\Omega$ is the set of all possible outcomes. An \textbf{event} $A \subseteq \Omega$ is a subset of outcomes. A probability function $P$ assigns $P(A) \in [0,1]$ with $P(\Omega) = 1$.
\end{rldefinition}
\pause
\begin{itemize}
  \item \textbf{Complement rule:} $P(A^c) = 1 - P(A)$.
  \pause
  \item \textbf{Addition rule:} $P(A \cup B) = P(A) + P(B) - P(A \cap B)$.
  \pause
  \item \textbf{Multiplication rule:} $P(A \cap B) = P(A)\,P(B \mid A)$.
\end{itemize}
\end{frame}

\begin{frame}{Mutual Exclusivity vs. Independence}
\begin{itemize}
  \item $A$ and $B$ are \textbf{mutually exclusive} if $A \cap B = \emptyset$, so $P(A \cap B) = 0$.
  \pause
  \item $A$ and $B$ are \textbf{independent} if $P(A \cap B) = P(A)\,P(B)$, equivalently $P(B \mid A) = P(B)$.
  \pause
  \item These are not the same: two mutually exclusive events with nonzero probability are always dependent, since $P(A\cap B) = 0 \neq P(A)P(B)$.
\end{itemize}
\pause
\begin{rlpitfall}
Confusing ``mutually exclusive'' with ``independent'' is a common error. If two events can never happen together, learning that one happened tells you the other definitely did not --- that is strong dependence, not independence.
\end{rlpitfall}
\end{frame}

\begin{frame}{Empirical vs. Classical Probability}
\begin{itemize}
  \item \textbf{Classical probability}: $P(A) = |A| / |\Omega|$ when all outcomes are equally likely (e.g. a fair die).
  \pause
  \item \textbf{Empirical probability}: $P(A) \approx \frac{\text{count of } A \text{ in observed data}}{\text{total observations}}$, estimated from a sample.
  \pause
  \item Machine learning almost always relies on empirical probability: class frequencies in a training set, token frequencies in a corpus, click-through rates in logged data.
\end{itemize}
\pause
\begin{rlquestion}
If a spam classifier is trained on a dataset where 90\% of emails are spam, is the empirical $P(\text{spam})$ estimated from that dataset a good estimate of $P(\text{spam})$ for real-world email traffic?
\end{rlquestion}
\end{frame}
